\documentclass[letterpaper, 12pt]{article}
% ============================================================
%  CONFIGURACION COMUN PARA SOLUCIONARIOS PEP 1
% ============================================================

\usepackage[utf8]{inputenc}
\usepackage[spanish]{babel}
\usepackage{amsmath, amssymb, amsfonts}
\usepackage{geometry}
\usepackage{fancyhdr}
\usepackage{enumitem}
\usepackage{graphicx}
\usepackage{hyperref}
\usepackage{needspace}
\usepackage{parskip}
\usepackage{xcolor}
\usepackage{tcolorbox}
\usepackage{eso-pic}
\usepackage{tikz}
\tcbuselibrary{skins, breakable}

\geometry{letterpaper, margin=1.7cm, top=3cm, bottom=2.5cm, headheight=2cm}

\definecolor{celesteSuave}{HTML}{F0F7FF}
\definecolor{azulFuerte}{HTML}{1A5276}
\definecolor{enlace}{HTML}{1A5276}

\newcommand{\R}{\mathbb{R}}
\newcommand{\Dom}{\operatorname{Dom}}
\newcommand{\Rec}{\operatorname{Rec}}
\newcommand{\Cod}{\operatorname{Cod}}

\newcommand{\logoup}{../../../../../template/images/logo_up.png}
\newcommand{\logousach}{../../../../../template/images/logo_usach.png}
\newcommand{\logofondo}{../../../../../template/images/logo_up.png}
\newcommand{\institucion}{Usach Premium}
\newcommand{\curso}{C\'alculo I}
\newcommand{\autor}{Jaime Riquelme}
\newcommand{\semestre}{1er. Semestre de 2026}
\newcommand{\fraseportada}{``Por y para estudiantes''}
\newcommand{\webinstitucion}{www.usachpremium.cl}
\newcommand{\instagraminstitucion}{https://www.instagram.com/usach.premium}
\newcommand{\transparencia}{0.05}

\hypersetup{
    colorlinks=true,
    linkcolor=enlace,
    urlcolor=enlace,
    citecolor=enlace,
    pdfborder={0 0 0},
}

\pagestyle{fancy}
\fancyhf{}
\fancyhead[L]{\includegraphics[height=1.2cm]{\logoup}}
\fancyhead[R]{%
    \begin{minipage}[b]{0.42\textwidth}
        \raggedleft
        \fontsize{9pt}{11pt}\selectfont
        \textbf{\color{azulFuerte}\institucion} \\
        \curso\ -- Solucionario PEP 1 \\
        \textit{\color{gray}@usach.premium}
    \end{minipage}\hspace{0.1cm}%
    \includegraphics[height=1.2cm]{\logousach}%
}
\fancyfoot[L]{\fontsize{9pt}{11pt}\selectfont \textit{\fraseportada}}
\fancyfoot[R]{\fontsize{9pt}{11pt}\selectfont P\'agina \thepage}
\renewcommand{\headrulewidth}{0.4pt}
\renewcommand{\footrulewidth}{0.4pt}

\fancypagestyle{plain}{
    \fancyhf{}
    \renewcommand{\headrulewidth}{0pt}
    \renewcommand{\footrulewidth}{0pt}
}

\newcommand{\MarcaAgua}{
    \AddToShipoutPicture{
        \AtPageCenter{
            \begin{tikzpicture}[remember picture, overlay]
                \node[opacity=\transparencia] at (0,0)
                    {\includegraphics[width=0.7\textwidth]{\logofondo}};
            \end{tikzpicture}
        }
    }
}

\newtcolorbox{solucionbox}[1]{
    colback=white,
    colframe=azulFuerte,
    coltitle=white,
    colbacktitle=azulFuerte,
    title=\textbf{#1},
    title after break=\textbf{#1} (continuaci\'on),
    fonttitle=\bfseries,
    arc=4pt,
    boxrule=1pt,
    enhanced,
    breakable,
    before={\par\Needspace{0.36\textheight}}
}

\newtcolorbox{respuestabox}{
    colback=celesteSuave,
    colframe=azulFuerte,
    boxrule=0.8pt,
    arc=4pt
}

\newcommand{\portadasolucionario}[2]{
    \thispagestyle{plain}
    \AddToShipoutPicture*{%
        \put(54, 700){\includegraphics[width=2cm]{\logoup}}
        \put(420, 706){\includegraphics[width=5cm]{\logousach}}
    }
    \vspace*{0.5cm}
    \begin{center}
        \color{azulFuerte}
        {\huge \textbf{SOLUCIONARIO}} \\[0.5cm]
        {\Huge \textbf{\curso}} \\[0.3cm]
        {\Large #1} \\[0.3cm]
        {\large #2} \\[1.4cm]
        \begin{tcolorbox}[colback=celesteSuave, colframe=azulFuerte, arc=10pt,
            width=0.82\textwidth, center, boxrule=1.5pt]
            \centering
            \Large\textbf{Desarrollo paso a paso}
        \end{tcolorbox}
    \end{center}
    \vspace{1.4cm}
    \begin{center}
        {\Large \textbf{\semestre}} \\[0.7cm]
        {\large \textbf{\autor\ -- \institucion}}
    \end{center}
    \vfill
    \begin{center}
        {\color{azulFuerte}\hrule height 0.5pt}
        \vspace{0.3cm}
        \begin{minipage}{0.45\textwidth}
            \centering
            \textbf{Instagram:} \\
            \small\href{\instagraminstitucion}{@usach.premium}
        \end{minipage}
        \begin{minipage}{0.45\textwidth}
            \centering
            \textbf{Web:} \\
            \small\href{http://\webinstitucion}{\webinstitucion}
        \end{minipage}
    \end{center}
    \newpage
    \MarcaAgua
}

\newcommand{\respuesta}[1]{
    \begin{respuestabox}
        \textbf{Respuesta:} #1
    \end{respuestabox}
}


\begin{document}
\portadasolucionario{Gu\'ia de Preparaci\'on PEP 1}{Secci\'on II -- Dominio, recorrido y composici\'on}

\begin{solucionbox}{Ejercicio 2.1}
Halle el dominio natural de
\[
    h(x)=\sqrt{\frac{x+4}{2-x}}.
\]

\textbf{Soluci\'on:}

Para que $h(x)$ est\'e bien definida en $\R$, el radicando debe ser mayor o igual que cero:
\[
    \frac{x+4}{2-x}\geq 0.
\]
Adem\'as, el denominador no puede anularse:
\[
    2-x\neq 0
    \quad\Longrightarrow\quad
    x\neq 2.
\]
Los puntos cr\'iticos son $x=-4$ y $x=2$.

\[
\begin{array}{c|ccc}
    & (-\infty,-4) & (-4,2) & (2,+\infty) \\ \hline
    x+4 & - & + & + \\
    2-x & + & + & - \\ \hline
    \dfrac{x+4}{2-x} & - & + & -
\end{array}
\]

Como necesitamos signo no negativo, se incluye $x=-4$ porque anula el numerador, pero no se incluye $x=2$ porque anula el denominador.

\respuesta{$\Dom(h)=[-4,2)$.}
\end{solucionbox}

\begin{solucionbox}{Ejercicio 2.2}
Considere $f(x)=\sqrt{5-x}$ y $g(x)=\sqrt{x^2-1}$. Determine $\Dom(f\circ g)$.

\textbf{Soluci\'on:}

Recordemos que
\[
    \Dom(f\circ g)=\{x\in\Dom(g):g(x)\in\Dom(f)\}.
\]

Primero, determinamos los dominios de $f$ y $g$.

Para $f$:
\[
    5-x\geq 0
    \quad\Longleftrightarrow\quad
    x\leq 5.
\]
Por lo tanto,
\[
    \Dom(f)=(-\infty,5].
\]

Para $g$:
\[
    x^2-1\geq 0
    \quad\Longleftrightarrow\quad
    (x-1)(x+1)\geq 0.
\]
Luego,
\[
    \Dom(g)=(-\infty,-1]\cup[1,+\infty).
\]

Ahora exigimos que $g(x)\in\Dom(f)$:
\[
    \sqrt{x^2-1}\leq 5.
\]
Como ambos lados son no negativos, elevamos al cuadrado:
\[
    x^2-1\leq 25
    \quad\Longleftrightarrow\quad
    x^2\leq 26
    \quad\Longleftrightarrow\quad
    -\sqrt{26}\leq x\leq \sqrt{26}.
\]

Finalmente,
\[
    \Dom(f\circ g)=
    \left((-\infty,-1]\cup[1,+\infty)\right)
    \cap[-\sqrt{26},\sqrt{26}].
\]

\respuesta{$\Dom(f\circ g)=[-\sqrt{26},-1]\cup[1,\sqrt{26}]$.}
\end{solucionbox}

\begin{solucionbox}{Ejercicio 2.3}
Considere
\[
    f(x)=\sqrt{\frac{x-1}{x+3}},
    \qquad
    g(x)=2x-5.
\]
Determine $\Dom(f\circ g)$, halle expl\'icitamente $(f\circ g)(x)$, calcule $(f\circ g)(4)$ y determine si existe una preimagen de $1$.

\textbf{Soluci\'on:}

Primero determinamos $\Dom(f)$. Para que $f$ est\'e definida:
\[
    \frac{x-1}{x+3}\geq 0,
    \qquad x\neq -3.
\]
Los puntos cr\'iticos son $x=-3$ y $x=1$, y se obtiene
\[
    \Dom(f)=(-\infty,-3)\cup[1,+\infty).
\]
Como $g(x)=2x-5$ es lineal,
\[
    \Dom(g)=\R.
\]

Por definici\'on,
\[
    \Dom(f\circ g)=\{x\in\R:2x-5\in(-\infty,-3)\cup[1,+\infty)\}.
\]
Entonces resolvemos:
\[
    2x-5<-3
    \quad\vee\quad
    2x-5\geq 1.
\]
De aqu\'i,
\[
    2x<2 \quad\vee\quad 2x\geq 6,
\]
es decir,
\[
    x<1 \quad\vee\quad x\geq 3.
\]
Por lo tanto,
\[
    \Dom(f\circ g)=(-\infty,1)\cup[3,+\infty).
\]

Ahora calculamos la regla:
\[
    (f\circ g)(x)=f(g(x))=f(2x-5).
\]
Luego,
\[
    (f\circ g)(x)
    =\sqrt{\frac{(2x-5)-1}{(2x-5)+3}}
    =\sqrt{\frac{2x-6}{2x-2}}
    =\sqrt{\frac{x-3}{x-1}}.
\]

Para la imagen de $4$:
\[
    (f\circ g)(4)
    =\sqrt{\frac{4-3}{4-1}}
    =\sqrt{\frac13}.
\]

Para la preimagen de $1$, buscamos $x\in\Dom(f\circ g)$ tal que
\[
    \sqrt{\frac{x-3}{x-1}}=1.
\]
Elevando al cuadrado:
\[
    \frac{x-3}{x-1}=1.
\]
Esto implica
\[
    x-3=x-1,
\]
lo cual es imposible. Por lo tanto, no existe preimagen de $1$.

\respuesta{$\Dom(f\circ g)=(-\infty,1)\cup[3,+\infty)$, $(f\circ g)(x)=\sqrt{\dfrac{x-3}{x-1}}$, $(f\circ g)(4)=\sqrt{\dfrac13}$ y no existe preimagen de $1$.}
\end{solucionbox}

\begin{solucionbox}{Ejercicio 2.4}
Sea
\[
    f(x)=1+\sqrt{\frac{x+1}{x-2}}.
\]
Determine $\Dom(f)$ y la preimagen de $3$, si existe.

\textbf{Soluci\'on:}

Para determinar el dominio, pedimos
\[
    \frac{x+1}{x-2}\geq 0,
    \qquad x\neq 2.
\]
Los puntos cr\'iticos son $x=-1$ y $x=2$.

\[
\begin{array}{c|ccc}
    & (-\infty,-1) & (-1,2) & (2,+\infty) \\ \hline
    x+1 & - & + & + \\
    x-2 & - & - & + \\ \hline
    \dfrac{x+1}{x-2} & + & - & +
\end{array}
\]

Por lo tanto,
\[
    \Dom(f)=(-\infty,-1]\cup(2,+\infty).
\]

Ahora buscamos la preimagen de $3$, es decir, resolvemos
\[
    f(x)=3.
\]
Entonces
\[
    1+\sqrt{\frac{x+1}{x-2}}=3
    \quad\Longleftrightarrow\quad
    \sqrt{\frac{x+1}{x-2}}=2.
\]
Elevando al cuadrado:
\[
    \frac{x+1}{x-2}=4.
\]
Luego,
\[
    x+1=4(x-2)
    \quad\Longleftrightarrow\quad
    x+1=4x-8.
\]
Por tanto,
\[
    9=3x
    \quad\Longleftrightarrow\quad
    x=3.
\]
Como $3\in\Dom(f)$, la preimagen existe.

\respuesta{$\Dom(f)=(-\infty,-1]\cup(2,+\infty)$ y la preimagen de $3$ es $x=3$.}
\end{solucionbox}

\begin{solucionbox}{Ejercicio 2.5}
Considere
\[
    p(x)=\sqrt{6-\sqrt{x+2}}.
\]
Determine $\Dom(p)$, indique si $p$ es creciente o decreciente en su dominio, y halle $\Rec(p)$.

\textbf{Soluci\'on:}

Primero, para que la ra\'iz interior est\'e definida:
\[
    x+2\geq 0
    \quad\Longleftrightarrow\quad
    x\geq -2.
\]
Adem\'as, para que la ra\'iz exterior est\'e definida:
\[
    6-\sqrt{x+2}\geq 0.
\]
Esto equivale a
\[
    6\geq \sqrt{x+2}.
\]
Como ambos lados son no negativos, elevamos al cuadrado:
\[
    36\geq x+2
    \quad\Longleftrightarrow\quad
    x\leq 34.
\]
Por lo tanto,
\[
    \Dom(p)=[-2,34].
\]

Para estudiar la monoton\'ia, sean $a,b\in[-2,34]$ tales que $a<b$. Entonces:
\[
    a<b
    \quad\Longrightarrow\quad
    a+2<b+2.
\]
Como la ra\'iz cuadrada es creciente:
\[
    \sqrt{a+2}<\sqrt{b+2}.
\]
Multiplicando por $-1$ y sumando $6$:
\[
    6-\sqrt{a+2}>6-\sqrt{b+2}.
\]
Aplicando nuevamente la ra\'iz cuadrada, que es creciente en $[0,+\infty)$:
\[
    \sqrt{6-\sqrt{a+2}}>
    \sqrt{6-\sqrt{b+2}}.
\]
Es decir,
\[
    p(a)>p(b).
\]
Por lo tanto, $p$ es decreciente en todo su dominio.

Como $p$ es decreciente,
\[
    \Rec(p)=[p(34),p(-2)].
\]
Calculamos:
\[
    p(-2)=\sqrt{6-\sqrt0}=\sqrt6,
\]
\[
    p(34)=\sqrt{6-\sqrt{36}}=\sqrt{6-6}=0.
\]
Por lo tanto,
\[
    \Rec(p)=[0,\sqrt6].
\]

\respuesta{$\Dom(p)=[-2,34]$, $p$ es decreciente y $\Rec(p)=[0,\sqrt6]$.}
\end{solucionbox}

\begin{solucionbox}{Ejercicio 2.6}
Sea
\[
    f(x)=\sqrt{\frac{4}{\sqrt{x-2}}-1},
    \qquad
    g(x)=x+3.
\]
Halle $\Dom(f)$, $\Rec(f)$ y $\Dom(f\circ g)$, indicando adem\'as la regla de asignaci\'on de $f\circ g$.

\textbf{Soluci\'on:}

Para que $f$ est\'e definida, primero necesitamos
\[
    \sqrt{x-2}
\]
en el denominador, por lo que
\[
    x-2>0
    \quad\Longleftrightarrow\quad
    x>2.
\]
Adem\'as, el radicando exterior debe ser mayor o igual que cero:
\[
    \frac{4}{\sqrt{x-2}}-1\geq 0.
\]
Entonces
\[
    \frac{4}{\sqrt{x-2}}\geq 1.
\]
Como $\sqrt{x-2}>0$, multiplicamos sin cambiar el sentido:
\[
    4\geq \sqrt{x-2}.
\]
Elevando al cuadrado:
\[
    16\geq x-2
    \quad\Longleftrightarrow\quad
    x\leq 18.
\]
Por lo tanto,
\[
    \Dom(f)=(2,18].
\]

Para hallar el recorrido, notemos que si $x\in(2,18]$, entonces
\[
    0<\sqrt{x-2}\leq 4.
\]
Luego,
\[
    \frac{4}{\sqrt{x-2}}\geq 1.
\]
Por lo tanto,
\[
    \frac{4}{\sqrt{x-2}}-1\geq 0.
\]
Cuando $x=18$, se obtiene $f(18)=0$. Cuando $x\to 2^+$, el denominador $\sqrt{x-2}\to0^+$ y la expresi\'on crece sin cota superior. De esta forma,
\[
    \Rec(f)=[0,+\infty).
\]

Ahora,
\[
    \Dom(f\circ g)=\{x\in\Dom(g):g(x)\in\Dom(f)\}.
\]
Como $g(x)=x+3$ tiene dominio $\R$:
\[
    x+3\in(2,18].
\]
Esto equivale a
\[
    2<x+3\leq 18
    \quad\Longleftrightarrow\quad
    -1<x\leq 15.
\]
Por lo tanto,
\[
    \Dom(f\circ g)=(-1,15].
\]
Finalmente,
\[
    (f\circ g)(x)=f(x+3)
    =\sqrt{\frac{4}{\sqrt{(x+3)-2}}-1}
    =\sqrt{\frac{4}{\sqrt{x+1}}-1}.
\]

\respuesta{$\Dom(f)=(2,18]$, $\Rec(f)=[0,+\infty)$, $\Dom(f\circ g)=(-1,15]$ y $(f\circ g)(x)=\sqrt{\dfrac{4}{\sqrt{x+1}}-1}$.}
\end{solucionbox}

\end{document}
